% SOURCE_AUTHORITY: sga5_fr_workpass.tex, lines 11028--11055 (LF-normalized unit).
% SOURCE_SHA256: 791F4EFFC5E02832D5D77ED03518C8156D6F07E4C8238B03545DB93D883FBB28
\emph{Demonstração.} Demostremos (9.1.1). Graças à fórmula de projeção para $p$, basta ver que $p_*(c_r(\check F))=1$. Para isso, basta reproduzir a demonstração de 8.4.1, usando as relações ([CH] 4 (17)). A igualdade entre o primeiro e o terceiro membro de (9.1.2) não é senão o lema 3 de [1]. Repitindo o argumento da demonstração de (8.4.1), vê-se que
\[
c_r(\check F)=\sum_{i=0}^r p^*(c_i(\check N))\xi^{r-i},
\tag*{(9.1.2) bis}
\]
o que termina a demonstração de (9.1.2). Aplicando o funtor $s^*$ aos dois membros extremos de (9.1.2), vê-se, pois $p\circ s=\operatorname{id}$, que (9.1.3) será consequência da igualdade
\[
s^*\left(\sum_{i=0}^{r-1}p^*(c_i(\check N))\xi^{r-i}\right)=0,
\]
e, com maior razão, das igualdades
\[
s^*(\xi^i)=0\qquad(i>0).
\tag{9.1.5}
\]
Demostremos (9.1.5). Denotando por $\alpha:V(N)\to\widehat N$ a inclusão canônica, é obviamente $\alpha^*(\mathcal O_P(1))\simeq\mathcal O_{V(N)}$, de onde
\[
\alpha^*(\xi)=\alpha^*(c_1(\mathcal O_P(1)))=c_1(\mathcal O_{V(N)})=0
\]
pela funtorialidade das classes de Chern. A afirmação resulta imediatamente da definição de $s$. Demonstremos (9.1.4). Para isso utilizaremos a igualdade
\[
\xi\,c_r(\check F)=c_{r+1}(p^*(\check N)\oplus\mathcal O_P)=0,
\tag{9.1.6}
\]
que resulta imediatamente da sequência exata (9.1.0). O elemento $z$ admite uma decomposição única da forma
\[
z=\sum_{i=0}^r p^*(a_i)\xi^i,\qquad a_i\in C^*(S),
\]
([CH] 4 Th. 1). De (9.1.5) resulta que $s^*(z)=a_0$. Por outro lado, a igualdade (9.1.6) mostra que $z\,c_r(\check F)=p^*(a_0)c_r(\check F)$, e conclui-se por (9.1.1).
